\section{Preguntas y respuestas: context finito, memoria externa y siguientes pasos}

La exposición principal termina alrededor de 1:00; las preguntas posteriores se concentran en tres límites: cómo distinguir el origen de los token multimodales, cómo ampliar un context finito y qué investigará Karpathy a continuación. Como algunas preguntas del público no se entienden, esta sección conserva solo los puntos didácticos que pueden reconstruirse con fiabilidad a partir de las respuestas, sin suponer las preguntas que faltan.

\subsection{El context no tiene que crecer sin límite}

Karpathy propone una idea que en ese momento era speculative: mantener finito el context length y enseñar al modelo a emitir instrucciones especiales de scratchpad para escribir en un notebook externo lo que debe recordar y consultarlo después. Así separa el «estado de corto plazo dentro de la cabeza» de una memory externa de lectura y escritura, y anticipa el diseño de interfaz de tool use y de los memory-augmented agent, aunque en clase no presentó ninguna implementación ni experimento.

\teachervoice{Su analogía es esta: una persona no guarda todos los hechos en la cabeza, sino que usa una libreta. El modelo también puede aprender a emitir marcas de inicio y fin (start/end) de scratchpad, y un sistema externo se encarga de guardar, indexar y volver a insertar el contenido. Lo importante es que el protocolo pueda supervisarse, direccionarse y verificarse; darle al modelo un prompt más largo no equivale a tener una memoria de largo plazo confiable.}

\subsection{Epistemic honesty: si no se sabe, se dice claramente}

Cuando le preguntan por los internals de ChatGPT, Karpathy responde directamente que no los conoce; cuando le preguntan por S4, tampoco finge estar familiarizado. Al final vuelve a nanoGPT, en el que trabaja: reescribir minGPT como un código más pequeño, más moderno y con el que sea más fácil reproducir GPT, y dice que le interesa la iteración continua de productos de lenguaje. Esta forma de responder fija el mismo criterio para estas notas: no completar detalles de un sistema más allá de la evidencia pública.

\lecturefigure{slide-61-thank-you.jpg}{Diapositiva final: una vez entendida la arquitectura, «go forth and transform».}{Video oficial de Stanford Online 01:00:31--01:00:45.}

La diapositiva final casi no tiene contenido técnico, pero cumple una función dentro del curso: pasar de una clase introductoria a las clases posteriores sobre cada área. Después de lo expuesto, ``transform'' ya no significa solo cambiar de backbone, sino completar un diseño de cuatro pasos: definir los elementos, inyectar la posición y el tipo, establecer la connectivity y elegir el objective; después, verificar que el enrutamiento de la attention, el cálculo de la MLP, la estabilidad de la optimización y el costo de hardware funcionen en conjunto.

\subsection{Resumen de la sección}

Las preguntas y respuestas convierten los límites de la arquitectura en interfaces de sistema: lo multimodal requiere type/position information, un context finito puede combinarse con memory externa y protocolos de herramientas, y las implementaciones internas desconocidas deben seguir tratándose como desconocidas. nanoGPT ofrece el camino ejecutable más corto para seguir aprendiendo: primero reproducir un modelo pequeño y luego poner a prueba con experimentos las intuiciones sobre attention, mask, context y scaling.
